\documentclass{article}
\usepackage[utf8]{inputenc}
\usepackage[margin=2cm,includefoot]{geometry}
\usepackage{booktabs}
\usepackage{graphicx}
\usepackage{amsmath, amssymb, amsthm}
\usepackage{physics}
\usepackage[shortlabels]{enumitem}
\usepackage{cancel}
\usepackage{lastpage}
\usepackage{braket}
\usepackage{mathrsfs}
\usepackage{bm}
\usepackage{subcaption}
\usepackage{float}

\usepackage{fancyhdr}
\pagestyle{fancy}
\newcommand{\AssignmentNum}{3}
\lhead{Stromingsleer en transportverschijnselen (NS-265B)\\
2025/2026}
\rhead{Tereza Bílková (7137257), Bela F. Brunner (8002193),\\ Niels Epema (2927578), Adam Zich (0074187)}
\fancyfoot{}
\fancyfoot[R]{Page \thepage /\pageref{LastPage}}

\newcommand{\vba}{\mathbf{A}}
\newcommand{\vbb}{\mathbf{B}}
\newcommand{\C}{\mathbf{C}}
\newcommand{\x}{\mathbf{\hat{x}}}
\newcommand{\y}{\mathbf{\hat{y}}}
\newcommand{\z}{\mathbf{\hat{z}}}
\newcommand{\vr}{\mathbf{\hat r}}
\newcommand{\thet}{\boldsymbol{{\hat \theta}}}
\newcommand{\ph}{\boldsymbol{{\hat \phi}}}
\newcommand{\ba}{\mathbf{a}}
\newcommand{\bb}{\mathbf{b}}
\newcommand{\vl}{\mathbf{l}}
\newcommand{\vv}{\mathbf{v}}

\newcommand{\R}{\mathbb{R}}
\newcommand{\p}{\mathbb{P}}
\newcommand{\N}{\mathbb{N}}
\newcommand{\Z}{\mathbb{Z}}
\newcommand{\Q}{\mathbb{Q}}
\newcommand{\I}{\mathbb{I}}
\newcommand{\E}{\mathbb{E}}
\renewcommand{\Pr}{\textbf{Pr}}
\newcommand{\Ra}{\textbf{Ra}}
\newcommand{\bracket}[1]{\left(#1\right)}

\newcommand{\mbeq}{\overset{!}{=}}
\newcommand{\PATH}{figures/}

\begin{document}
    \setlength{\headheight}{23.10004pt}
    \begin{center}
    \LARGE{\bfseries{ Project \AssignmentNum}}
    \line(1,0){430}
    \end{center}

\section{Analytic part}

\subsection*{a)}
\subsection*{b)}
    To find the expression for $\frac{d\Psi}{dt}$ we start of with the equation from problem (a)
    \begin{gather}
        \frac{1}{\Pr}(\partial_t(\nabla^2 \psi) + u \partial_x(\nabla^2 \psi) + w \partial_z(\nabla^2 \psi) = \partial_x^2(\nabla^2 \psi) + \partial_z^2 (\nabla^2 \psi) + \Ra \partial_x \vartheta \\
    \end{gather}
    Using the expression for $\psi(x,z,t)$ and $\vartheta(x,z,t)$
    \begin{gather*}
        \psi(x,z,t) = \Psi(t) \sin(kx) \sin(\pi z) \\
        \vartheta(x,z,t) = \Theta_1(t) \cos(kx) \sin(\pi z) - \Theta_2(t) \sin(2\pi z) 
    \end{gather*}
    
    we can calculate the components of the first equation seperately
    \begin{align*}
        \nabla^2 \psi &= -\Psi(t) k^2 \sin(kx) \sin(\pi z) - \Psi(t) \pi^2 \sin(k x) \sin(\pi z) \\
        %
        \partial_t(\nabla^2 \psi) &= \partial_t (-\Psi(t) k^2 \sin(kx) \sin(\pi z) - \Psi(t) \pi^2 \sin(k x) \sin(\pi z)) \\
        &= -\frac{d\Psi(t)}{dt} \sin(k x) \sin(\pi z) (k^2 + \pi^2) \\
        %
        \partial_x(\nabla^2 \psi) &= \partial_x (-\Psi(t) k^2 \sin(kx) \sin(\pi z) - \Psi(t) \pi^2 \sin(k x) \sin(\pi z)) \\
        &= \Psi(t) \cos(k x) \sin(\pi z) (k^3 + k\pi^2) \\
        %
        \partial_z (\nabla^2 \psi) &= \partial_z (-\Psi(t) k^2 \sin(kx) \sin(\pi z) - \Psi(t) \pi^2 \sin(k x) \sin(\pi z)) \\
        &= \Psi(t) \sin(k x) \cos(\pi z) (k^2 \pi + \pi^3) \\
        %
        \partial_x^2 (\nabla^2 \psi) &= \partial_x^2 (-\Psi(t) k^2 \sin(kx) \sin(\pi z) - \Psi(t) \pi^2 \sin(k x) \sin(\pi z)) \\
        &= \Psi(t) \sin(k x) \sin(\pi z) (k^2 \pi^2 + \pi^4) \\
        %
        \partial_z^2 (\nabla^2 \psi) &= \partial_z^2 (-\Psi(t) k^2 \sin(kx) \sin(\pi z) - \Psi(t) \pi^2 \sin(k x) \sin(\pi z)) \\
        &= \Psi(t) \sin(k x) \sin(\pi z) (k^4 + k^2 \pi^2) \\
        %
        \Ra \partial_x \vartheta &= \Ra \partial_x (\Theta_1(t) \cos(kx) \sin(\pi z) - \Theta_2(t) \sin(2\pi z)) \\
        &=-\Ra (\Theta_1(t) k \sin(kx) \sin(\pi z)) \\
        %
        u &= -\partial_z \psi = -\Psi(t) \pi \sin(kx) \cos(\pi z) \\
        %
        w &= \Psi(t) k \cos(kx) \sin(\pi z)
    \end{align*}
    Combining this into the equation, while multiplying everything by $\Pr$ and taking everything to the left side gives us 
    \begin{align*}
        &-\frac{d\Psi(t)}{dt} \sin(k x) \sin(\pi z) (k^2 + \pi^2) + u \Psi(t) \cos(k x) \sin(\pi z) (k^3 + k\pi^2) + w \Psi(t) \sin(k x) \cos(\pi z) (k^2 \pi + \pi^3) \\ 
        &\qquad - \Pr \Psi(t) \sin(k x) \sin(\pi z) (k^2 \pi^2 + \pi^4) - \Pr \Psi(t) \sin(k x) \sin(\pi z) (k^4 + k^2 \pi^2) + \Pr \Ra (\Theta_1(t) k \sin(kx) \sin(\pi z)) = 0
    \end{align*}
    Now we can take out $\Psi(t)$ and divide by the term with the $\frac{d\Psi(t)}{dt}$
    \begin{align*}
        &-\frac{d\Psi(t)}{dt} \sin(k x) \sin(\pi z) (k^2 + \pi^2) + \Psi(t) (u \cos(k x) \sin(\pi z) (k^3 + k\pi^2) + w \sin(k x) \cos(\pi z) (k^2 \pi + \pi^3) \\
        &\qquad - \Pr \sin(k x) \sin(\pi z) (k^2 \pi^2 + \pi^4) - \Pr \sin(k x) \sin(\pi z) (k^4 + k^2 \pi^2)) + \Pr \Ra (\Theta_1(t) k \sin(kx) \sin(\pi z)) = 0 \\
        %
        &-\frac{d\Psi(t)}{dt} + \Psi(t) (u \frac{\cos(k x)}{\sin(k x)} (\frac{k^3 + k\pi^2}{k^2 + \pi^2}) + w \frac{\cos(\pi z)}{\sin(\pi z)} (\frac{k^2 \pi + \pi^3}{k^2 + \pi^2}) - \Pr (\frac{k^2 \pi^2 + \pi^4}{k^2 + \pi^2}) - \Pr (\frac{k^4 + k^2 \pi^2}{k^2 + \pi^2})) \\
        & \qquad+ \Pr \frac{\Ra \Theta_1(t) k}{k^2 + \pi^2} = 0
    \end{align*}
    Substitution in the expressions for $u$ and $w$, taking $\frac{d\Psi(t)}{dt}$ to the other side and simplifying the equation gives us
    \begin{align*}
        \frac{d\Psi(t)}{dt} &= \Psi(t) (-\Psi(t) \pi \sin(kx) \cos(\pi z) \frac{\cos(k x)}{\sin(k x)} (\frac{k^3 + k\pi^2}{k^2 + \pi^2}) + \Psi(t) k \cos(kx) \sin(\pi z) \frac{\cos(\pi z)}{\sin(\pi z)} (\frac{k^2 \pi + \pi^3}{k^2 + \pi^2}) \\
        &\qquad- \Pr (\frac{K^4 + 2k^2 \pi^2 + \pi^4}{k^2 + \pi^2}) + \Pr \frac{\Ra \Theta_1(t) k}{k^2 + \pi^2} \\
        %
        \frac{d\Psi(t)}{dt} &= \Psi(t) (-\Psi(t) \cos(\pi z) \cos(k x) (\frac{k^3 \pi + k\pi^3}{k^2 + \pi^2}) + \Psi(t) \cos(kx) \cos(\pi z) (\frac{k^3 \pi + k \pi^3}{k^2 + \pi^2}) - \Pr (\frac{k^4 + 2k^2 \pi^2 + \pi^4}{k^2 + \pi^2})) \\
        &\qquad + \Pr \frac{\Ra \Theta_1(t) k}{k^2 + \pi^2} \\
        %
        \frac{d\Psi(t)}{dt} &= \Psi(t) (\Psi(t) \cos(\pi z) \cos(k x) (\frac{-k^3 \pi - k\pi^3 + k^3 \pi + k \pi^3}{k^2 + \pi^2}) - \Pr (\frac{(k^2 + \pi^2)^2}{k^2 + \pi^2})) + \Pr \frac{\Ra \Theta_1(t) k}{k^2 + \pi^2} \\
        %
        \frac{d\Psi(t)}{dt} &= - \Pr (k^2 + \pi^2) \Psi(t) + \Pr \frac{k \Ra}{k^2 + \pi^2} \Theta_1(t) 
    \end{align*}

We want to show that the equations for $\Theta_1(t)$ and $\Theta_2(t)$ are given by:
\begin{equation}\label{T1}
    \frac{d\Theta_1}{dt} = k\Psi - (k^2 + \pi^2)\Theta_1 - \pi k \Psi \Theta_2
\end{equation}
\begin{equation}\label{T2}
    \frac{d\Theta_2}{dt} = \frac{1}{2}k\pi \Psi \Theta_1 - 4\pi^2\Theta_2
\end{equation}
To do this, we plug the expression for temperature
\begin{equation}
    \vartheta(x,z,t) = \Theta_1(t)\cos kx \sin \pi z - \Theta_2(t)\sin 2\pi z
\end{equation}
into
\begin{equation}\label{temp}
    \partial_t \vartheta + u\partial_x \vartheta + w\partial_z \vartheta - w = \partial^2_x \vartheta + \partial^2_z \vartheta
\end{equation} 
Computing the derivatives:
\begin{align*}
    \partial_t \vartheta &= \frac{d\Theta_1(t)}{dt}\cos kx \sin \pi z - \frac{d\Theta_2}{dt}\sin 2\pi z \\
    \partial_x \vartheta &= -k \Theta_1 \sin kx \sin \pi z \\
    \partial_z \vartheta &= \pi \Theta_1 \cos kx \cos \pi z - 2\pi \Theta_2 \cos 2\pi z \\
    \partial^2_x \vartheta &= -k^2 \Theta_1 \cos kx \sin \pi z \\
    \partial^2_z \vartheta &= -\pi^2 \Theta_1 \cos kx \sin \pi z + 4\pi^2 \Theta_2 \sin 2\pi z
\end{align*}
Using the definitions of u and w:
\begin{align*}
    u = -\partial_z \psi = -\Psi(t) \pi \sin kx \cos \pi z \\
    w = \partial_x \psi = \Psi(t) k \cos kx \sin \pi z
\end{align*}
We start off by plugging these derivatives, along with the definitions of u and w into expression \ref{temp}. \\
Doing this, the left-hand-side becomes:
\begin{align*}
    \left(\frac{d\Theta_1(t)}{dt}\cos kx \sin \pi z - \frac{d\Theta_2}{dt}\sin 2\pi z\right) + \left(-\Psi(t) \pi \sin kx \cos \pi z\right)\left(-k \Theta_1 \sin kx \sin \pi z\right) \\ + \left(\Psi(t) k \cos kx \sin \pi z\right)\left(\pi \Theta_1 \cos kx \cos \pi z - 2\pi \Theta_2 \cos 2\pi z\right) - \left(\Psi(t) k \cos kx \sin \pi z\right)
\end{align*}
and the right-hand-side becomes:
\begin{align*}
    \left(-k^2 \Theta_1 \cos kx \sin \pi z\right) + \left(-\pi^2 \Theta_1 \cos kx \sin \pi z + 4\pi^2 \Theta_2 \sin 2\pi z\right)
\end{align*}
We can use the indentity $\sin 2 \pi z = 2\cos \pi z \sin \pi z$ to simplify the terms $\sin 2\pi z$ on both sides, which allows us to divide both sides by $\sin \pi z$. Doing this and multiplying out terms in brackets gives:
Left-hand-side:
\begin{align*}
    \frac{d\Theta_1}{dt}\cos kx - 2\frac{d\Theta_2}{dt}\cos \pi z + k\pi \Psi\Theta_1 \sin^2 kx \cos \pi z + k\pi\Psi\Theta_1 \cos^2 kx \cos \pi z - 2k\pi \Psi\Theta_2 \cos kx \cos 2\pi z - k\Psi \cos kx
\end{align*}
We can merge 
\begin{align*}
    -k \pi \Psi\Theta_1  \sin^2 kx \cos \pi z - k \pi \Psi \Theta_1  \cos^2 kx \cos \pi z
\end{align*}
into 
\begin{align*}
    -k \pi \Psi \Theta_1 \cos \pi z
\end{align*}
which simplifies our expression to
\begin{align*}
    \frac{d\Theta_1}{dt}\cos kx - 2\frac{d\Theta_2}{dt}\cos \pi z -k \pi \Psi \Theta_1 \cos \pi z - 2k\pi \Psi\Theta_2 \cos kx \cos 2\pi z - k\Psi \cos kx
\end{align*}
Right-hand-side:
\begin{align*}
    -k^2 \Theta_1 \cos kx - \pi^2 \Theta_1 \cos kx + 8\pi^2 \Theta_2 \cos \pi z
\end{align*}
where $-k^2\Theta_1 \cos kx - \pi^2 \Theta_1 \cos kx = -\Theta_1\cos kx \left(k^2 + \pi^2\right)$, so this becomes:
\begin{align*}
    -\Theta_1\cos kx \left(k^2 + \pi^2\right) + 8\pi^2 \Theta_2 \cos \pi z
\end{align*}
Setting the two sides equal and equating all the coefficients in front of $\cos kx$ and $\cos \pi z$, respectively, yields the following pair of equalities:
\begin{equation*}
    \frac{d\Theta_1}{dt} = -\Theta_1 \left(k^2 + \pi^2\right) + 2\pi k \Psi \Theta_2 \cos 2\pi z + k\Psi \\
\end{equation*}
\begin{equation*}
        \frac{d\Theta_2}{dt} = -\frac{k}{2} \pi \Psi \Theta_1 - 4\pi^2 \Theta_2 
\end{equation*}
The second equation we obtain matches \ref{T2}, however the first doesn't yet match \ref{T1}. To make them match, we use the identity to reexpress the second term
\begin{align*}
    \cos 2\pi z \sin \pi z = \frac{1}{2}\left(\sin 3 \pi z - \sin \pi z\right)
\end{align*}
This is done as follows:
\begin{align*}
    2 \pi k \Psi \Theta_2 \cos 2\pi z = 2 \pi k \Psi \Theta_2 \cos 2\pi z \frac{\sin \pi z}{\sin \pi z}
\end{align*}
using the identity, this becomes
\begin{align*}
    2 \pi k \Psi \Theta_2 \cos 2\pi z \frac{\sin \pi z}{\sin \pi z} = \pi k \Theta_2 \Psi \frac{\sin 3\pi z - \sin \pi z}{\sin \pi z}
\end{align*}
since we are only considering the first 2 Fourier modes in the z-direction, we can drop the term $\sin 3 \pi z$, and we get:
\begin{align*}
    \pi k \Theta_2 \Psi \frac{\sin 3\pi z - \sin \pi z}{\sin \pi z} = \pi k \Theta_2 \Psi
\end{align*}
This means if we now substitute this in place of the second term, we get equations which match \ref{T1} and \ref{T2}:
\begin{equation*}
    \frac{d\Theta_1}{dt} = -\Theta_1 \left(k^2 + \pi^2\right) + \pi k \Theta_2 \Psi + k\Psi \\
\end{equation*}
\begin{equation*}
        \frac{d\Theta_2}{dt} = -\frac{k}{2} \pi \Psi \Theta_1 - 4\pi^2 \Theta_2 
\end{equation*}

\subsection*{c)}


\section{Computational part}


\end{document}
